% ============================================================
%  共源 NMOS 放大器 SPICE 仿真器设计报告
%  编译方式: xelatex report.tex  (需运行两次以生成目录)
% ============================================================
\documentclass[12pt, a4paper]{article}

% ── 中英文字体 ──────────────────────────────────────────────
\usepackage[UTF8, fontset=none]{ctex}
\usepackage{fontspec}
\setmainfont{Times New Roman}
\setCJKmainfont[BoldFont=STHeiti, ItalicFont=STKaiti]{STFangsong}   % 仿宋（正文）
\setCJKsansfont[BoldFont=STHeiti]{STHeiti}                           % 黑体（标题）
\setCJKmonofont[BoldFont=STHeiti]{STFangsong}

% ── 页面版式 ────────────────────────────────────────────────
\usepackage[top=2.5cm, bottom=2.5cm, left=3cm, right=3cm, headheight=15pt]{geometry}
\usepackage{setspace}
\setstretch{1.35}

% ── 数学 ────────────────────────────────────────────────────
\usepackage{amsmath, amssymb, bm}

% ── 图表 ────────────────────────────────────────────────────
\usepackage{graphicx}
\usepackage{float}
\usepackage{caption}
\usepackage{subcaption}
\usepackage{booktabs}
\usepackage{array}
\captionsetup{font=small, labelfont=bf, skip=6pt}

% ── 电路图 ──────────────────────────────────────────────────
\usepackage{tikz}
\usepackage[european, straightvoltages]{circuitikz}
\usetikzlibrary{arrows.meta, positioning, calc}

% ── 代码 ────────────────────────────────────────────────────
\usepackage{listings}
\usepackage{xcolor}
\lstset{
  basicstyle=\small\ttfamily,
  keywordstyle=\color{blue!70},
  commentstyle=\color{green!50!black},
  stringstyle=\color{orange!80!black},
  frame=single, rulecolor=\color{gray!40},
  breaklines=true, numbers=left,
  numberstyle=\tiny\color{gray},
  xleftmargin=1.5em, xrightmargin=0.5em,
  showstringspaces=false, tabsize=4
}

% ── 复旦 Logo ────────────────────────────────────────────────
\usepackage{fdulogo}

% ── 单位 ────────────────────────────────────────────────────
\usepackage{siunitx}
\sisetup{mode=math, per-mode=symbol}

% ── 超链接 ──────────────────────────────────────────────────
\usepackage[colorlinks, linkcolor=black, citecolor=black,
            urlcolor=blue!70!black]{hyperref}

% ── 章节标题格式（编号用 Times New Roman，标题文字用黑体）──────
\usepackage{titlesec}
\titleformat{\section}
  {\large\bfseries\sffamily}
  {{\normalfont\large\bfseries\thesection}}{0.8em}{}
\titleformat{\subsection}
  {\normalsize\bfseries\sffamily}
  {{\normalfont\normalsize\bfseries\thesubsection}}{0.8em}{}
\titleformat{\subsubsection}
  {\normalsize\bfseries\sffamily}
  {{\normalfont\normalsize\bfseries\thesubsubsection}}{0.8em}{}
\titlespacing{\section}{0pt}{1.2ex plus 0.5ex}{0.6ex}
\titlespacing{\subsection}{0pt}{0.9ex plus 0.3ex}{0.4ex}

% ── 目录标题字号 ─────────────────────────────────────────────
\usepackage{tocloft}
\renewcommand{\cfttoctitlefont}{\Large\bfseries\sffamily}
\setlength{\cftaftertoctitleskip}{1em}

% ── 页脚（仅页码，无页眉）────────────────────────────────────
\pagestyle{plain}

% ── 自定义命令 ──────────────────────────────────────────────
\newcommand{\vth}{V_{\mathrm{th}}}
\newcommand{\vgs}{V_{\mathrm{GS}}}
\newcommand{\vds}{V_{\mathrm{DS}}}
\newcommand{\vsb}{V_{\mathrm{SB}}}
\newcommand{\vov}{V_{\mathrm{ov}}}
\newcommand{\id}{I_{\mathrm{D}}}
\newcommand{\gm}{g_{\mathrm{m}}}
\newcommand{\gds}{g_{\mathrm{ds}}}
\newcommand{\gmb}{g_{\mathrm{mb}}}
\newcommand{\vT}{V_{\mathrm{T}}}

% ============================================================
\begin{document}
% ============================================================

% ── 封面 ────────────────────────────────────────────────────
\begin{titlepage}
  \centering
  \vspace*{1.5cm}

  % 校徽 + 校名横排（x/y 为路径坐标单位，max coord ≈1039，0.07pt ≈ 2.5cm）
  \fduemblem+[color=FudanBlue, x=0.07pt, y=0.07pt]
  \hspace{0.5cm}
  \fduname+[color=FudanBlue, x=0.07pt, y=0.07pt]

  \vspace{1.0cm}
  {\normalsize\sffamily 集成电路与微纳电子创新学院\par}

  \vspace{1.8cm}
  \rule{0.62\linewidth}{0.5pt}
  \vspace{0.6cm}

  {\LARGE\bfseries\sffamily 共源 NMOS 放大器\\[0.5em]
   SPICE 仿真器设计报告\par}

  \vspace{0.6cm}
  \rule{0.62\linewidth}{0.5pt}
  \vspace{1.6cm}

  {\large\sffamily 器件模型与 SPICE 仿真 \quad 2026 春\par}

  \vspace{2.0cm}
  {\normalsize
  \begin{tabular}{rl}
    小组成员 & 王耀冬（23300760013）\quad 陈星（23300760005）\quad 韩卓骏（23300760003）\\[0.5em]
    代码仓库 & \url{https://github.com/LArielOzjH/nano-spice} \\[0.5em]
    仿真语言 & Python 3 \;(NumPy + Matplotlib) \\
  \end{tabular}
  \par}

  \vfill
\end{titlepage}

% ── 目录 ────────────────────────────────────────────────────
\tableofcontents
\clearpage

% ============================================================
\section{引言}
% ============================================================

本报告记录了一个从零构建的共源 NMOS 放大器 SPICE 仿真器的完整设计过程。
仿真器以改进节点分析法（Modified Nodal Analysis，MNA）为核心，
配合 Newton-Raphson 迭代求解非线性 DC 方程组，
并在此基础上扩展了二阶物理效应模型、结二极管模型
以及复数域 AC 小信号分析。

\textbf{主要目标：}
\begin{itemize}
  \item 建立可收敛的非线性 DC 工作点求解框架，覆盖截止区、线性区、饱和区；
  \item 逐步引入亚阈值导通、迁移率衰减、DIBL、体效应、寄生串联电阻、结二极管等物理效应，
        每项均可独立开关，形成清晰的效应对比；
  \item 实现基于电荷分区的电容模型，并将 MNA 推广至复数域，完成 AC 频率扫描及带宽计算。
\end{itemize}

% ============================================================
\section{电路拓扑与节点定义}
% ============================================================

\subsection{电路结构}

仿真目标为图~\ref{fig:circuit} 所示共源 NMOS 放大器。
负载电阻 $R_D = 10\,\si{\kilo\ohm}$ 接于 $V_{DD}=5\,\mathrm{V}$ 与输出节点 VOUT 之间；
NMOS 晶体管 M0 的栅极由输入电压 $V_{IN}$ 直接驱动，源极接地。

\begin{figure}[H]
\centering
\begin{circuitikz}[american, scale=1.1, font=\small,
                   every node/.style={font=\small}]

  % ── 主支路：VDD → RD → VOUT → M0 → GND ─────────────
  %    MOSFET 中心固定在 (4, 3)
  \node[nmos] (M) at (4, 3) {};

  % RD：从漏极向上画，直接到 VDD 标注
  \draw (M.drain)
        to[R, l_={$R_D{=}10\,\si{\kilo\ohm}$}] ++(0, 3.2)
        node[above]{$V_{DD}{=}5\,\mathrm{V}$};

  % VOUT 节点：在漏极处打点，向右引出标注
  \filldraw (M.drain) circle (1.6pt);
  \draw (M.drain) -- ++(2.4, 0) node[right]{VOUT\;($V_2$)};

  % 源极接地
  \draw (M.source) -- ++(0, -0.7) node[ground]{};

  % ── VIN 支路：位于左侧，与漏极对齐的高度之下 ─────────
  %    利用垂直坐标对齐：(x |- node) 取 node 的 y 坐标
  % nmos 中心在 (4,3)，gate 与中心同高 → y=3
  \draw (1.5, 0.6) node[ground]{}
        to[V, l=$V_{IN}$]
        (1.5, 3) -- (M.gate);

  % ── 标注 ─────────────────────────────────────────────
  % M0：放在晶体管右侧中部（漏极在上，此处不与 VOUT 引线重叠）
  \node[right=5pt] at (M.east) {M0};

\end{circuitikz}
\caption{共源 NMOS 放大器电路图（$V_{th}=0.8\,\mathrm{V}$，
$k_p=100\,\mu\mathrm{A/V^2}$，$W/L=10$，$\lambda=0.02\,\mathrm{V^{-1}}$）。
引入寄生电阻时，在 VOUT 与漏极间插入 $R_d$（本征节点 $V_5$），
在源极与地间插入 $R_s$（本征节点 $V_4$）。}
\label{fig:circuit}
\end{figure}

\subsection{节点与未知量定义}

为统一处理有无寄生电阻两种情况，MNA 系统始终维持 7 个未知量：
\[
  \bm{x} = \bigl[V_1,\; V_2,\; V_3,\; V_4\,(S'),\; V_5\,(D'),\; I_{V_{IN}},\; I_{V_{DD}}\bigr]^T
\]
各节点含义如表~\ref{tab:nodes} 所示。当 $R_s=0$ 时令 $V_4=0$（约束方程），
当 $R_d=0$ 时令 $V_5=V_2$，系统退化为等效的无寄生情形。

\begin{table}[H]
\centering
\caption{MNA 未知量定义}
\label{tab:nodes}
\begin{tabular}{cll}
\toprule
索引 & 物理含义 & 备注 \\
\midrule
$V_1$ & Gate 电压 & 由 $V_{IN}$ 电压源约束 \\
$V_2$ & VOUT / 外部漏极 & 输出节点 \\
$V_3$ & $V_{DD}$ & 由 $V_{DD}$ 电压源约束 \\
$V_4$ & 本征源极 $S'$ & $R_s=0$ 时约束为 $0$ \\
$V_5$ & 本征漏极 $D'$ & $R_d=0$ 时约束为 $V_2$ \\
$I_{V_{IN}}$ & 流过 $V_{IN}$ 源的电流 & KVL 引入的辅助变量 \\
$I_{V_{DD}}$ & 流过 $V_{DD}$ 源的电流 & KVL 引入的辅助变量 \\
\bottomrule
\end{tabular}
\end{table}

% ============================================================
\section{改进节点分析法（MNA）}
% ============================================================

\subsection{基本原理}

MNA 将 KCL（节点电流守恒）和 KVL（电压源约束）联立，
对线性网络可直接建立 $\bm{A}\bm{x}=\bm{b}$；
对含 MOSFET 等非线性元件的电路，则在每次 Newton-Raphson 迭代时
用伴随模型（小信号线性化）替换非线性器件，重建线性方程组。

\subsection{七变量方程组}

以节点 $V_4$（$S'$）和 $V_5$（$D'$）的本征管脚为基础，
NMOS 的内部电压为：
\[
  V_{GS,\mathrm{int}} = V_1 - V_4,\quad
  V_{DS,\mathrm{int}} = V_5 - V_4,\quad
  V_{SB} = V_4
\]
各行方程对应含义如下（以 $R_s>0,\;R_d>0$ 的一般情况为例）：

\begin{align}
  \text{行 0 (节点1 KCL):}&\quad I_{V_{IN}} = 0 \label{eq:row0}\\
  \text{行 1 (节点2 KCL):}&\quad (V_3-V_2)G_{load} - (V_2-V_5)G_{rd} = 0 \label{eq:row1}\\
  \text{行 2 (节点3 KCL):}&\quad -(V_3-V_2)G_{load} + I_{V_{DD}} = 0 \label{eq:row2}\\
  \text{行 3 (节点4 KCL):}&\quad I_D + I_{BS} - V_4 G_{rs} = 0 \label{eq:row3}\\
  \text{行 4 (节点5 KCL):}&\quad (V_2-V_5)G_{rd} - I_D + I_{BD} = 0 \label{eq:row4}\\
  \text{行 5 (KVL\ }V_{IN}\text{):}&\quad V_1 = V_{IN} \label{eq:row5}\\
  \text{行 6 (KVL\ }V_{DD}\text{):}&\quad V_3 = V_{DD} \label{eq:row6}
\end{align}

其中 $G_{load}=1/R_D$，$G_{rs}=1/R_s$，$G_{rd}=1/R_d$；
$I_{BD}$、$I_{BS}$ 为 BD/BS 结二极管电流（第~\ref{sec:diode} 节详述）。

\subsection{伴随模型与线性化雅可比}

在工作点 $\bm{x}^{(k)}$ 处，NMOS 漏极电流线性化为：
\[
  I_D \approx I_{eq} + \gm V_1 + (-\gm - \gds + \gmb)V_4 + \gds V_5
\]
其中伴随电流源常数为：
\[
  I_{eq} = I_D^{(k)} - \gm^{(k)} V_{GS,\mathrm{int}}^{(k)}
           - \gds^{(k)} V_{DS,\mathrm{int}}^{(k)} - \gmb^{(k)} V_{SB}^{(k)}
\]
将此线性化关系代入方程组，即得每步迭代的线性方程 $\bm{A}^{(k)}\bm{x}^{\mathrm{new}}=\bm{b}^{(k)}$，
用标准 LU 分解求解。

% ============================================================
\section{NMOS 晶体管模型}
% ============================================================

\subsection{基础模型（Shichman-Hodges + CLM）}
\label{sec:base}

基础漏极电流公式为：
\[
  \id = k_p \frac{W}{L}
        \left(V_{gsteff}\,V_{ds,\mathrm{eff}} - \tfrac{1}{2}V_{ds,\mathrm{eff}}^2\right)
        (1 + \lambda V_{DS})
\]
其中 $\lambda$ 为沟道长度调制系数，$(1+\lambda V_{DS})$ 项描述沟道夹断后电流随
$V_{DS}$ 的缓慢增大（有限输出电阻 $r_o = 1/(\lambda I_D)$）。

器件参数取值：$V_{th}=0.8\,\mathrm{V}$，$k_p=100\,\mu\mathrm{A/V^2}$，
$W/L=10$，$\lambda=0.02\,\mathrm{V^{-1}}$。

\subsection{softmin 平滑技术}
\label{sec:softmin}

传统 SPICE 模型在线性区与饱和区边界存在一阶导数不连续（$\partial I_D/\partial V_{DS}$ 跳变），
导致 Newton-Raphson 在该区域收敛缓慢甚至振荡。
本仿真器采用 \textit{softmin} 函数替代硬截断：
\[
  V_{ds,\mathrm{eff}} = \frac{1}{2}\!\left(
    V_{DS} + V_{gsteff}
    - \sqrt{(V_{DS} - V_{gsteff})^2 + \delta^2}
  \right)
\]
其中 $\delta=0.02\,\mathrm{V}$ 为平滑参数。
当 $V_{DS} \ll V_{gsteff}$ 时 $V_{ds,\mathrm{eff}}\approx V_{DS}$（线性区），
当 $V_{DS} \gg V_{gsteff}$ 时 $V_{ds,\mathrm{eff}}\approx V_{gsteff}$（饱和区），
过渡区内函数光滑可微，雅可比矩阵连续，保证 Newton-Raphson 的二次收敛性。

偏导数解析形式为：
\[
  \frac{\partial V_{ds,\mathrm{eff}}}{\partial V_{DS}} = \frac{1}{2}\!\left(
    1 - \frac{V_{DS}-V_{gsteff}}{\sqrt{(V_{DS}-V_{gsteff})^2+\delta^2}}
  \right),\quad
  \frac{\partial V_{ds,\mathrm{eff}}}{\partial V_{gsteff}} = 1 -
  \frac{\partial V_{ds,\mathrm{eff}}}{\partial V_{DS}}
\]

\subsection{亚阈值导通}
\label{sec:subthreshold}

标准 Shichman-Hodges 模型在 $V_{GS}<V_{th}$ 时令 $I_D=0$，忽略了弱反型区的指数型电流。
引入统一有效栅压 $V_{gsteff}$，使弱反型与强反型连续过渡：
\[
  V_{gsteff} = n \vT \ln\!\left(1 + e^{V_{ov}/(n\vT)}\right)
\]
其中 $V_{ov}=V_{GS}-V_{th}$，$n$ 为亚阈值斜率因子（典型值 $1.1\sim2.0$），
$\vT=kT/q\approx25.85\,\mathrm{mV}$（室温）。

当 $V_{ov}\gg n\vT$ 时 $V_{gsteff}\approx V_{ov}$（强反型），
当 $V_{ov}\ll0$ 时 $V_{gsteff}\approx n\vT e^{V_{ov}/(n\vT)}$（弱反型指数律），
两端均无突变，微分连续。对应的 $\gm$ 修正因子为 sigmoid 函数：
\[
  \frac{\partial V_{gsteff}}{\partial V_{GS}} = \frac{e^{V_{ov}/(n\vT)}}{1+e^{V_{ov}/(n\vT)}}
  \in (0,\,1)
\]

\subsection{迁移率衰减}

高横向电场（大 $V_{GS}$）导致载流子表面散射增强，有效迁移率下降：
\[
  \mu_{eff} = \frac{\mu_0}{1 + \theta \cdot V_{ov}},\quad
  k_{p,\mathrm{eff}} = \frac{k_p}{1 + \theta \cdot V_{ov}},\quad V_{ov}>0
\]
$\theta$ 为迁移率衰减系数（典型值 $0.05\sim0.5\,\mathrm{V^{-1}}$）。
对 $\gm$ 的影响：跨导在大 $V_{GS}$ 下不再线性增大，而是趋于饱和，
这与实际测量一致。

\subsection{漏致势垒降低（DIBL）}

短沟道器件中，漏极电场穿透沟道影响源端势垒，导致阈值电压随 $V_{DS}$ 下降：
\[
  V_{th,\mathrm{eff}} = V_{th0} - \eta \cdot V_{DS}
\]
$\eta$ 为 DIBL 系数（典型值 $0.01\sim0.1\,\mathrm{V/V}$）。
与沟道长度调制（$\lambda$ 项）不同，DIBL 对阈值电压的改变影响
$V_{ov}$ 的计算，从而放大了饱和区电流对 $V_{DS}$ 的敏感度，
给出更真实的输出电导。
DIBL 对 $\gds$ 的贡献为：
\[
  \gds\big|_{\mathrm{DIBL}} = \gm \cdot \eta
\]

\subsection{体效应}

源极与体（衬底）之间的偏置电压 $V_{SB}$ 使耗尽层变宽，阈值电压升高：
\[
  V_{th,\mathrm{eff}} = V_{th0} + \gamma\left(\sqrt{|V_{SB}|+2\phi_F} - \sqrt{2\phi_F}\right)
\]
$\gamma$ 为体效应系数（典型值 $0.3\sim1.0\,\mathrm{V^{1/2}}$），$2\phi_F$ 为强反型所需表面势。
相应的体效应跨导为：
\[
  \gmb = \gm \cdot \frac{\gamma}{2\sqrt{|V_{SB}|+2\phi_F}}
\]
在本电路（源极接地，$V_{SB}=V_4$）中，当引入寄生源极电阻 $R_s$ 后
$V_4>0$，体效应自然产生；无 $R_s$ 时 $V_4=0$，体效应为零。

\subsection{寄生串联电阻（Rs/Rd）}

实际工艺中，源极与漏极的欧姆接触引入串联电阻 $R_s$、$R_d$。
处理方式是在 MNA 中显式引入两个本征节点 $S'$（$V_4$）和 $D'$（$V_5$），
以线性电阻连接至外部节点（图~\ref{fig:circuit}），MOSFET 模型仅操作内部电压。
这一做法无需修改晶体管模型本身，并且天然产生源极体效应（$V_4>0$）和有效跨导衰减：
\[
  g_{m,\mathrm{eff}} \approx \frac{\gm}{1 + \gm R_s}
\]

\subsection{BD/BS 结二极管}
\label{sec:diode}

漏极/源极 n\textsuperscript{+} 扩散区与 p 型衬底形成反偏 PN 结。
本仿真器采用理想二极管模型：
\[
  I_{BD} = I_S\!\left(e^{V_{BD}/\vT} - 1\right),\quad
  V_{BD} = V_B - V_{D'} = -V_5
\]
采用伴随模型（Norton 等效）线性化后，在本征漏极节点的 MNA stamp 为：
\begin{align}
  A[D',D'] &\mathrel{-}= G_{BD} = \frac{I_S}{\vT}e^{V_{BD}/\vT} \\
  b[D'] &\mathrel{-}= I_{BD,0} = I_{BD} - G_{BD}\cdot V_{BD}
\end{align}
BS 结处理方式相同（$V_{BS}=-V_4$，stamp 至节点 $S'$）。
对指数项作 $40\vT$ 限幅以防止正偏时数值溢出。
当 $I_S=0$ 时所有二极管路径自动关闭，与基础模型完全兼容。

% ============================================================
\section{Newton-Raphson 迭代求解器}
% ============================================================

\subsection{算法框架}

设非线性方程组为 $\bm{F}(\bm{x})=\bm{0}$，Newton-Raphson 迭代为：
\[
  \bm{x}^{(k+1)} = \bm{x}^{(k)} - \bm{J}^{-1}\!\left(\bm{x}^{(k)}\right)\bm{F}\!\left(\bm{x}^{(k)}\right)
\]
其中雅可比矩阵 $\bm{J}=\partial\bm{F}/\partial\bm{x}$ 即为 MNA 矩阵 $\bm{A}$，
由 MOSFET 的解析导数（$\gm,\gds,\gmb$）及各电阻的电导填充，无需数值差分。
解析雅可比使迭代在工作点附近具有二次收敛速度，通常仅需 2\;--\;6 步即可收敛。

\subsection{阻尼策略}

初始猜测偏离工作点较远时，全步 Newton 更新可能使残差增大。
仿真器采用残差驱动的 $\lambda$-halving 策略：若完整步长使
$\|\bm{F}(\bm{x}^{(k+1)})\| > \|\bm{F}(\bm{x}^{(k)})\|$，
则将步长折半 $\lambda \leftarrow \lambda/2$，
最多折半 10 次，最终取满足条件的最大 $\lambda$。

\subsection{SPICE 标准收敛判据}

采用与商业 SPICE 一致的逐分量收敛判据（课程第~11 章 p.~12）：
\[
  |\Delta x_i| < \varepsilon_{a,i} + \varepsilon_r \cdot
  \min\!\left(|x_i^{(k)}|,\; |x_i^{(k+1)}|\right), \quad \forall i
\]
其中：
\begin{itemize}
  \item $\varepsilon_{a}$：绝对容差，电压变量取 $10^{-6}\,\mathrm{V}$，电流变量取 $10^{-12}\,\mathrm{A}$；
  \item $\varepsilon_r = 10^{-3}$：相对容差（0.1\%），与 HSPICE/PSpice 默认值一致；
  \item $\min(\cdot)$ 而非 $\max(\cdot)$：避免因一端接近零而过于宽松。
\end{itemize}
所有分量同时满足则判定收敛，最大迭代次数设为 100。

% ============================================================
\section{电容模型与 AC 小信号分析}
% ============================================================

\subsection{基于电荷分区的栅极电容模型}

根据课程第~4 章的电荷守恒分区方法，各工作区内本征栅极电容值如下：

\begin{table}[H]
\centering
\caption{各工作区栅极本征电容分配}
\label{tab:caps}
\begin{tabular}{lccc}
\toprule
工作区 & $C_{gs,\mathrm{int}}$ & $C_{gd,\mathrm{int}}$ & $C_{gb}$ \\
\midrule
截止区 ($V_{ov}<0$) & $0$ & $0$ & $C_{gg}$ \\
线性区 ($V_{DS}<V_{ov}$) & $C_{gg}/2$ & $C_{gg}/2$ & $0$ \\
饱和区 ($V_{DS}>V_{ov}$) & $\tfrac{2}{3}C_{gg}$ & $0$ & $0$ \\
\bottomrule
\end{tabular}
\end{table}

为保持数值连续性，区域切换采用 sigmoid 函数平滑：
\[
  C_{gs} = C_{gg}\,\sigma_{on}\!\left(\tfrac{1}{2}+\tfrac{1}{6}\sigma_{sat}\right) + C_{gso},\quad
  C_{gd} = C_{gg}\,\sigma_{on}\cdot\tfrac{1}{2}(1-\sigma_{sat}) + C_{gdo},\quad
  C_{gb} = C_{gg}(1-\sigma_{on})
\]
其中 $\sigma_{on}=\sigma(V_{ov}/w)$，$\sigma_{sat}=\sigma\bigl((V_{DS}-V_{DS,\mathrm{sat}})/w\bigr)$，
$w=50\,\mathrm{mV}$ 为过渡宽度；$C_{gso}$、$C_{gdo}$ 为工艺覆盖电容（overlap capacitance）。

\subsection{复数域 MNA}

将 DC 工作点处的小信号电导矩阵 $\bm{G}$（由 $\gm,\gds,G_{load}$等构成）
与电容矩阵 $\bm{C}$（由 $C_{gs},C_{gd},C_{gb}$ 的 stamp 构成）叠加，
得到导纳矩阵：
\[
  \bm{Y}(j\omega) = \bm{G} + j\omega\bm{C}
\]
对各频率点求解：
\[
  \bm{Y}(j\omega)\,\bm{x}_{ac} = \bm{b}_{ac}
\]
激励向量 $\bm{b}_{ac}$ 在 $V_{IN}$ 对应行置 $1$（单位 AC 激励），其余为 $0$。
解向量的 $V_2$ 分量即为输出电压，传递函数为：
\[
  H(j\omega) = \frac{V_2(\omega)}{1}
\]

$\bm{C}$ 矩阵中 $C_{gs}$、$C_{gd}$、$C_{gb}$ 的 stamp 规则（以 $R_d>0$ 为例）：
\begin{itemize}
  \item $C_{gs}$：连接 Gate ($V_1$) 与 $S'$ ($V_4$)，按差分电容 stamp（$[1,1],[4,4]$ 加，$[1,4],[4,1]$ 减）；
  \item $C_{gd}$：连接 Gate ($V_1$) 与 $D'$ ($V_5$)，stamp 类似；
  \item $C_{gb}$：Gate ($V_1$) 对地，仅 $[1,1]$ 加 $C_{gb}$。
\end{itemize}

\subsection{频率响应与 $-3\,\mathrm{dB}$ 带宽}

在 DC 工作点 $V_{IN}=1.5\,\mathrm{V}$（饱和区）处，以对数间隔取
$10^6\sim10^{12}\,\mathrm{Hz}$ 的频率点，逐点求解复数 MNA，计算：
\[
  |H(j\omega)|_{dB} = 20\log_{10}|H(j\omega)|
\]
在幅度曲线上通过线性插值定位 $-3\,\mathrm{dB}$ 频率点 $f_{-3\mathrm{dB}}$。

\subsection{单位电流增益截止频率 $f_T$}

$f_T$ 定义为短路电流增益降至 1 时的频率，解析式为：
\[
  f_T = \frac{\gm}{2\pi(C_{gs}+C_{gd})}
\]
该指标反映晶体管本征速度上限，不受外部负载影响。

% ============================================================
\section{仿真结果}
% ============================================================

\subsection{DC 工作点扫描}

图~\ref{fig:vout} 展示了 VOUT 与 VIN 的完整转移特性。
电路工作区间可明显分为三段：
$V_{IN}<0.8\,\mathrm{V}$ 为截止区（VOUT $\approx V_{DD}$）；
$0.8\,\mathrm{V}<V_{IN}<1.4\,\mathrm{V}$ 为饱和区（高增益线性放大）；
$V_{IN}>1.4\,\mathrm{V}$ 时 MOSFET 进入线性区（VOUT 被拉至低电位）。

\begin{figure}[H]
  \centering
  \includegraphics[width=0.82\textwidth]{../results/vout_vs_vin.png}
  \caption{VOUT vs.\ VIN 转移特性曲线（稀疏扫描点 + 连续曲线叠加）}
  \label{fig:vout}
\end{figure}

图~\ref{fig:smallsig} 给出了各 VIN 点的小信号参数。
跨导 $\gm$ 在饱和区中部（约 $V_{IN}=1.2\,\mathrm{V}$）达到峰值，
此后受迁移率衰减或进入线性区而下降；
增益 $|A_v|=\gm(r_o\|R_D)$ 的峰值约为 $16\,\mathrm{dB}$。

\begin{figure}[H]
  \centering
  \includegraphics[width=0.82\textwidth]{../results/smallsig_vs_vin.png}
  \caption{小信号参数 $g_m$、$r_o$、$|A_v|$ 随 VIN 的变化}
  \label{fig:smallsig}
\end{figure}

图~\ref{fig:newton} 记录了每个 VIN 点处 Newton-Raphson 的迭代次数。
绝大多数点在 3\;--\;4 次内收敛，过渡区边界处略有增加但不超过 8 次，
验证了 softmin + 解析雅可比 + SPICE 收敛判据的有效性。

\begin{figure}[H]
  \centering
  \includegraphics[width=0.72\textwidth]{../results/newton_iters.png}
  \caption{Newton-Raphson 迭代次数分布}
  \label{fig:newton}
\end{figure}

\subsection{二阶物理效应分析}

图~\ref{fig:effects_vout} 将基础模型与各二阶效应逐一对比。
各效应的定性影响总结如下：

\begin{itemize}
  \item \textbf{亚阈值导通}：在 $V_{IN}<V_{th}$ 区间产生指数型漏电流，
        VOUT 低于基础模型的 $V_{DD}$，转移特性更平滑；
  \item \textbf{迁移率衰减}：大 $V_{GS}$ 时电流受限，VOUT 的下降斜率变小，
        导致相同 VIN 下 VOUT 更高；
  \item \textbf{DIBL}：降低有效阈值，使转移特性向左平移，
        等效增大了饱和区的 $\gds$；
  \item \textbf{体效应}：仅在引入 $R_s$ 后 $V_4>0$ 时生效，升高 $V_{th}$，
        减小电流；
  \item \textbf{寄生串联电阻}：插入 $R_s$/$R_d$ 后等效跨导下降，
        放大器增益降低，VOUT 下降幅度变小。
\end{itemize}

\begin{figure}[H]
  \centering
  \includegraphics[width=0.95\textwidth]{../results/secondary_effects_vout.png}
  \caption{基础模型与各二阶物理效应对 VOUT 的影响对比}
  \label{fig:effects_vout}
\end{figure}

\begin{figure}[H]
  \centering
  \includegraphics[width=0.95\textwidth]{../results/secondary_effects_gm_id.png}
  \caption{各效应对 $g_m$ 和 $I_D$ 的影响对比}
  \label{fig:effects_gm}
\end{figure}

\subsection{AC 频率响应}

图~\ref{fig:bode} 为 AC Bode 图，工作点取 $V_{IN}=1.5\,\mathrm{V}$（饱和区）。

\begin{figure}[H]
  \centering
  \includegraphics[width=0.9\textwidth]{../results/ac_bode.png}
  \caption{AC Bode 图（幅度 + 相位），标注 $-3\,\mathrm{dB}$ 带宽与 $f_T$}
  \label{fig:bode}
\end{figure}

从图中可以读出：
低频增益约 $20\,\mathrm{dB}$（与 DC 小信号增益 $|A_v|=\gm(r_o\|R_D)\approx10$ 一致），
幅度在高频处单极点滚降；
$-3\,\mathrm{dB}$ 带宽 $f_{-3\mathrm{dB}}$ 由输出节点对地等效电容和负载电阻决定，
主要受 $C_{gd}$（Miller 等效）影响。
$f_T = \gm/(2\pi(C_{gs}+C_{gd}))$ 反映晶体管本征速度，
在所选工作点下约为数 GHz 量级。

图~\ref{fig:caps} 给出了栅极电容随 VIN 的变化，验证了第~\ref{sec:subthreshold} 节
中各区间的分配规律：截止区 $C_{gb}$ 占主导；进入饱和区后 $C_{gs}$ 约为 $2/3\,C_{gg}$，
$C_{gd}$ 趋近零（仅剩覆盖电容 $C_{gdo}$）。

\begin{figure}[H]
  \centering
  \includegraphics[width=0.82\textwidth]{../results/ac_caps_vs_vin.png}
  \caption{$C_{gs}$、$C_{gd}$、$C_{gb}$ 随 VIN 的变化（含覆盖电容）}
  \label{fig:caps}
\end{figure}

% ============================================================
\section{与 PSpice 仿真对比}
% ============================================================

\textit{（本节待补充。将对照 PSpice 仿真数据，在同一坐标系下比较 VOUT、$I_D$、$g_m$ 等曲线，
  并对误差来源进行定量分析。）}

% ============================================================
\section{结论}
% ============================================================

本仿真器从零实现了一个功能完整的共源 NMOS 放大器 SPICE 级仿真器，主要成果如下：

\begin{enumerate}
  \item \textbf{MNA 框架}：建立了 7 变量扩展 MNA，统一处理有无寄生电阻两种情况，
        结构清晰、可扩展性强；
  \item \textbf{物理模型}：在 Shichman-Hodges 基础上逐步引入了亚阈值导通、迁移率衰减、
        DIBL、体效应、寄生串联电阻、BD/BS 结二极管共 6 项二阶效应，
        每项均可独立开关并有精确解析导数；
  \item \textbf{求解器}：Newton-Raphson 结合 softmin 连续化、阻尼策略和 SPICE 标准收敛判据，
        收敛稳健，典型点仅需 3\;--\;4 次迭代；
  \item \textbf{AC 分析}：实现了基于电荷分区电容模型的复数域 MNA，
        可计算全频段 Bode 图、$-3\,\mathrm{dB}$ 带宽及 $f_T$。
\end{enumerate}

仿真代码已开源于 GitHub，结构模块化，便于后续扩展瞬态仿真（backward Euler）、
噪声分析及多晶体管电路。

% ============================================================
\section{小组分工}
% ============================================================

\begin{table}[H]
\centering
\caption{小组成员分工}
\begin{tabular}{lll}
\toprule
姓名 & 学号 & 主要负责内容 \\
\midrule
王耀冬 & 23300760013 &
  基础 MOSFET 模型（Shichman-Hodges + CLM）、softmin 平滑、\\
  & & MNA 矩阵构建、Newton-Raphson 求解器设计与调试 \\[0.4em]
陈星   & 23300760005 &
  二阶物理效应建模（亚阈值导通、迁移率衰减、DIBL、体效应）、\\
  & & BD/BS 结二极管伴随模型、SPICE 标准收敛判据实现 \\[0.4em]
韩卓骏 & 23300760003 &
  寄生串联电阻 $R_s/R_d$ 扩展、AC 小信号分析（电容模型、\\
  & & 复数域 MNA、Bode 图、$f_T$ 计算）、仿真脚本与报告撰写 \\
\bottomrule
\end{tabular}
\end{table}

% ============================================================
\clearpage
\section*{附录：代码仓库}
\addcontentsline{toc}{section}{附录：代码仓库}
% ============================================================

\begin{center}
  \url{https://github.com/LArielOzjH/nano-spice}
\end{center}

\medskip
\noindent 仓库目录结构：
\begin{lstlisting}[language=bash, caption={项目目录结构}]
nano-spice/
+-- simulate.py           # 主脚本：运行所有仿真并保存图形
+-- spice_sim/
|   +-- models/
|   |   +-- nmos.py       # NMOS 模型（基础 + 二阶效应 + 电容）
|   |   +-- resistor.py   # 线性电阻 MNA stamp
|   +-- core/
|   |   +-- mna.py        # MNA 矩阵构建（7 变量，含寄生电阻和二极管）
|   |   +-- newton.py     # Newton-Raphson 求解器（SPICE 收敛判据）
|   +-- analysis/
|       +-- dc_op.py      # 单点 DC 工作点
|       +-- dc_sweep.py   # DC 参数扫描
|       +-- ac_sweep.py   # AC 小信号频率扫描
+-- results/              # 自动生成的仿真图形（git-ignored）
\end{lstlisting}

\noindent 运行仿真：
\begin{lstlisting}[language=bash, caption={运行方式}]
pip install numpy matplotlib
python simulate.py     # 生成 results/ 下所有图形
\end{lstlisting}

\end{document}
